\documentclass[../../../include/open-logic-section]{subfiles}

\begin{document}

\olfileid{his}{set}{hilbertcurve}

\olsection{Lampiran: Kurva Hilbert kang Ngebaki Bidang}

Ing bab \olref[pathology]{sec}, wis kita sebutaké yèn buktiné Cantor
ngenani garis lan kothak kang cacah titiké padha persis
(\olref[his][set][cantorplane]{thm:cantorplane}) njalari Peano
takon apa ana \emph{kurva} kang ngebaki bidang. Ing taun 1890
dhèwèké olèh wangsulan iya. Ing pérangan iki kita njlèntrèhaké,
kanthi cara kang durung gumathok banget, carané mbangun kurva Hilbert
kang ngebaki bidang (kanthi sithik owah-owahan).\footnote{Kanggo
panjlèntrèhan kang luwih ketat, delengen \cite{Rose2010}.
Owah-owahan iku arupa tambahan pérangan abang ing kurva-kurva
ing ngisor iki. Kanthi mangkono, kontinuitas kurva rada luwih
gampang dipriksa.}

Kita bakal netepaké fungsi $h$ minangka limit runtutan fungsi
$h_1$, $h_2$, $h_3$, \dots\@ Dhisik kita njlèntrèhaké pambanguné.
Sabanjuré kita tuduhaké yèn fungsi iki ngebaki bidang, banjur
yèn fungsi iki pancèn sawijining kurva.

Daerah asil $h$ kita tetepaké kothak satuan, $\unitsquare$.
Iki panyedhakan kapisan marang $h$, yaiku $h_1$:
\begin{center}
	\begin{tikzpicture}
	\draw[gray] (0pt, 0pt) rectangle (80pt, 80pt);
	\draw[step = 40pt, gray, very thin] (0pt, 0pt) grid (80pt, 80pt);
	\draw[oldiagcolorC, thick] (20pt, 0)--(20pt, 20pt);
	\draw[oldiagcolorC, thick] (60pt, 0)--(60pt, 20pt);		
	\begin{scope}[xshift=20pt, yshift=20pt]
	\draw [black, thick, l-system={Hilbert curve, step=40pt, angle=90, axiom=L, order=1}]  lindenmayer system;
	\end{scope}
	\end{tikzpicture}
\end{center}
Supaya gampang ngetutaké, kothak iki kita bagi nganggo kisi
$2 \times 2$. Kurvané bisa kita anggep diwiwiti saka seprapat
kiwa ngisor, banjur menyang kiwa ndhuwur, tengen ndhuwur, lan
pungkasané tengen ngisor. Iki tahap pambangunan kapindho, yaiku $h_2$:
\begin{center}
	\begin{tikzpicture}
	\draw[gray] (0pt, 0pt) rectangle (80pt, 80pt);
	\draw[step = 20pt, gray, very thin] (0pt, 0pt) grid (80pt, 80pt);
	\draw[oldiagcolorC, thick] (0pt, 10pt)--(10pt, 10pt);
	\draw[oldiagcolorC, thick] (70pt, 10pt)--(80pt, 10pt);
	\draw[thick, oldiagcolorE] (10pt, 30pt)--(10pt, 50pt); 
	\draw[thick, oldiagcolorE] (30pt, 50pt)--(50pt, 50pt); 
	\draw[thick, oldiagcolorE] (70pt, 50pt)--(70pt, 30pt); \begin{scope}[xshift=10pt, yshift=30pt]
	\draw [thick, rotate=270,black, l-system={Hilbert curve, step=20pt, angle=90, axiom=L, order=1}]  lindenmayer system;
	\end{scope}
	\begin{scope}[xshift=10pt, yshift=50pt]
	\draw [thick, black, l-system={Hilbert curve, step=20pt, angle=90, axiom=L, order=1}]  lindenmayer system;
	\end{scope}
	\begin{scope}[xshift=50pt, yshift=50pt]
	\draw [thick, black, l-system={Hilbert curve, step=20pt, angle=90, axiom=L, order=1}]  lindenmayer system;
	\end{scope}
	\begin{scope}[xshift=70pt, yshift=10pt]
	\draw [thick, rotate=90,black, l-system={Hilbert curve, step=20pt, angle=90, axiom=L, order=1}]  lindenmayer system;
	\end{scope}
	\end{tikzpicture}
\end{center}
Werna-werna kang béda mbiyantu njlèntrèhaké carané mbangun $h_2$.
Dhisik, salinan pérangan $h_1$ kang dudu !!{colorC}, kanthi ukuran
dicilikaké, kita pasang ing kiwa ngisor, kiwa ndhuwur, tengen
ndhuwur, lan tengen ngisor kothak (digambar ireng). Banjur
patang gambar iku kita gandhèng nganggo garis !!{colorE}.
Pungkasané, gambar iki kita sambungaké marang wates kothak
nganggo garis !!{colorC}.

Saiki delengen $h_3$. Kaya déné $h_2$ dumadi saka patang salinan
pérangan $h_1$ kang dudu abang, kang dicilikaké lan digandhèng,
$h_3$ uga dumadi saka patang salinan pérangan $h_2$ kang dudu abang,
kanthi ukuran dicilikaké (digambar ireng). Salinan-salinan iku
banjur digandhèng nganggo garis !!{colorE}, lan pungkasané
disambungaké marang wates kothak nganggo garis !!{colorC}.
\begin{center}
	\begin{tikzpicture}
	\draw[gray] (0pt, 0pt) rectangle (80pt, 80pt);
	\draw[step = 10pt, gray, very thin] (0pt, 0pt) grid (80pt, 80pt);
	\draw[thick, oldiagcolorC](5pt, 0pt)--(5pt, 5pt); 
	\draw[thick, oldiagcolorC] (75pt, 0pt)--(75pt, 5pt); 
	\draw[thick, oldiagcolorE] (75pt, 45pt)--(75pt, 35pt);
	\draw[thick, oldiagcolorE] (5pt, 35pt)--(5pt, 45pt); 
	\draw[thick, oldiagcolorE] (35pt, 45pt)--(45pt, 45pt); 
	\draw[thick, oldiagcolorE] (75pt, 45pt)--(75pt, 35pt); \begin{scope}[xshift=5pt, yshift=35pt]
	\draw [rotate=270, thick, black, l-system={Hilbert curve, step=10pt, angle=90, axiom=L, order=2}]  lindenmayer system;
	\end{scope}
	\begin{scope}[xshift=5pt, yshift=45pt]
	\draw [thick, black, l-system={Hilbert curve, step=10pt, angle=90, axiom=L, order=2}]  lindenmayer system;
	\end{scope}
	\begin{scope}[xshift=45pt, yshift=45pt]
	\draw [thick, black, l-system={Hilbert curve, step=10pt, angle=90, axiom=L, order=2}]  lindenmayer system;
	\end{scope}
	\begin{scope}[xshift=75pt, yshift=5pt]
	\draw [thick, rotate=90,black, l-system={Hilbert curve, step=10pt, angle=90, axiom=L, order=2}]  lindenmayer system;
	\end{scope}
	\end{tikzpicture}
\end{center}
Saiki katon pola umum kanggo netepaké $h_{n+1}$ saka $h_n$.
Pungkasané, kurva $h$ \emph{dhewe} kita tetepaké kanthi nimbang
limit titik demi titik saka fungsi-fungsi $h_1$, $h_2$,
\dots\@ Tegesé, kanggo saben $x \in \unitline$:
\begin{align*}
	h(x) &= \lim_{n \rightarrow \infty} h_n(x)
\end{align*} 
Saiki kita jlèntrèhaké gagasan yèn kurva iki ngebaki bidang.
Nalika nggambar kurva $h_n$, kita masang kisi $2^n \times 2^n$
ing $\unitsquare$. Miturut téoréma Pythagoras, dawa diagonal
saben petak kisi yaiku:
\[
\sqrt{\left(\nicefrac{1}{2^{n}}\right)^2+\left(\nicefrac{1}{2^{n}}\right)^2} = 2^{(\frac{1}{2}-n)}
\]
lan cetha yèn $h_n$ ngliwati saben petak kisi. Mula, saben titik
ing $\unitsquare$ jaraké \emph{ora ngluwihi}
$2^{(\frac{1}{2}-n)}$ saka sawatara titik ing $h_n$. Saiki, $h$ wis
ditetepaké minangka limit saka fungsi-fungsi $h_1$, $h_2$,
$h_3$, \dots\@ Sketsa iki banjur nganggep yèn wates paling gedhé
jarak sawijining titik saka $h$ diwènèhaké déning:
\[
\lim_{n \rightarrow \infty} 2^{(\frac{1}{2}-n)} = 0.
\]
Tegesé, saben titik ing $\unitsquare$ dianggep nduwèni jarak
$0$ saka~$h$. Kanthi intuisi iki, saben titik ing $\unitsquare$
ana \emph{ing} kurva. Mula $h$ ngebaki bidang!{}

Isih perlu dituduhaké yèn $h$ pancèn sawijining \emph{kurva}.
Kanggo kuwi, kita kudu netepaké pangertèné. Definisi modhèren
adhedhasar definisi kang diwènèhaké Jordan ing taun 1887
(mung sawetara taun sadurungé kurva kapisan kang ngebaki bidang
diwènèhaké):

\begin{defn}
Kurva yaiku pemetaan kontinu saka $\unitline$ menyang $\Real^2$.
\end{defn}

Definisi iki cukup gampang dibayangaké: sawijining kurva iku
pemetaan kang ``lancar'' saka garis baku menyang bidang $\Real^2$.
Fungsi $h$ pancèn pemetaan saka $\unitline$ menyang $\Real^2$.
Supaya dadi kurva, kita isih kudu mbuktèkaké yèn $h$ kontinu.
Kontinuitas wis kita tetepaké ing \olref[limits]{sec} nganggo
notasi $\epsilon$/$\delta$. Katrangan ing ngisor iki mung
nyedhiyakaké gambaran bab carané kurva nyedhaki saben titik ing
kothak: \emph{Yèn sawijining titik $p$ ing $\unitsquare$ lan
tingkat katlitèn $\epsilon$ ditemtokaké, kita bisa nemu pérangan
mbukak saka $\unitline$ saéngga kanggo saben $x$ ing pérangan
iku, $h(x)$ ana ing sajroning jarak $\epsilon$ saka $p$.}

Dadi, anggepen $p$ lan $\epsilon$ wis ditemtokaké. Iki padha
karo nggambar bunderan kang puseuré $p$ lan radiusé $\epsilon$
ing $\unitsquare$. (Bunderan iki bisa ngliwati wates
$\unitsquare$, nanging kuwi ora dadi prakara.) Élinga yèn nalika
njlèntrèhaké fungsi $h_n$, kita nggambar kisi $2^n \times 2^n$
ing $\unitsquare$. Sepira waé ciliking $\epsilon$, ana $n$
kang cukup gedhé saéngga sawijining petak ing kisi
$2^n \times 2^n$ ing $\unitsquare$ iku mapan sakabèhé ing bunderan kang puseuré
$p$ lan radiusé $\epsilon$.

Dadi, pilih $n$ iku, lan anggep $I$ pérangan mbukak paling gedhé
saka $\unitline$ kang sakabèhé dipetakké déning $h_n$ menyang
petak kisi kang dimaksud. (Ing sketsa iki, anané interval
$(a,b)$ dimangertèni saka kasunyatan yèn $h_n$ ngliwati
saben petak ing kisi $2^n\times 2^n$.) Saiki kudu dituduhaké
yèn kanggo saben $x \in I$, titik $h(x)$ tetep ana ing petak
kisi kang padha. Kanggo \emph{iki}, cukup yèn $h_m(x)$ tetep
ana ing petak iku kanggo saben $m > n$. Yèn kita deleng
$h_m$ kanggo $m > n$, miturut pola gambar,
``pérangan kang padha'' saka interval satuan tetep dipetakké
menyang petak kisi kang padha, mung jaluré saya dawa lan
mlengkung. Pangrembug iki nerangaké intuisi pambangunané,
nanging durung dadi bukti kontinuitas fungsi iki kanthi syarat
épsilon-dèlta ing saben titik domain.

\end{document}
